\documentclass[11pt]{article}

\usepackage[T1]{fontenc}
\usepackage[utf8]{inputenc}
\usepackage{libertinus}
\usepackage[a4paper,margin=25mm,headheight=15pt]{geometry}
\usepackage{microtype}
\usepackage{amsmath,amssymb,mathtools,amsthm,bm}
\usepackage{mathrsfs}
\usepackage{booktabs,longtable,tabularx,array}
\usepackage{enumitem}
\usepackage{xcolor}
\usepackage[most]{tcolorbox}
\usepackage{listings}
\usepackage{fancyhdr}
\usepackage{hyperref}
\usepackage[nameinlink,noabbrev]{cleveref}
\usepackage{url}

\definecolor{Ink}{HTML}{1E293B}
\definecolor{Accent}{HTML}{1E5A7A}
\definecolor{AccentLight}{HTML}{EAF2F6}
\definecolor{RuleGray}{HTML}{CBD5E1}
\definecolor{CodeBg}{HTML}{F7F8FA}
\definecolor{CodeComment}{HTML}{4B766E}
\definecolor{CodeKeyword}{HTML}{784B84}

\hypersetup{
  colorlinks=true,
  linkcolor=Accent,
  citecolor=Accent,
  urlcolor=Accent,
  pdftitle={Asymptotic Transseries for the Inverses of the Gamma and Barnes G-Functions},
  pdfauthor={Research draft prepared for the ProveIt project}
}

\pagestyle{fancy}
\fancyhf{}
\fancyhead[L]{\small Inverse Gamma and Barnes--\(G\) transseries}
\fancyhead[R]{\small \thepage}
\renewcommand{\headrulewidth}{0.4pt}

\setlength{\parindent}{1.25em}
\setlength{\parskip}{0.15em}
\setlist{itemsep=0.25em,topsep=0.4em}
\allowdisplaybreaks

\newtheorem{theorem}{Theorem}[section]
\newtheorem{proposition}[theorem]{Proposition}
\newtheorem{lemma}[theorem]{Lemma}
\newtheorem{corollary}[theorem]{Corollary}
\theoremstyle{definition}
\newtheorem{definition}[theorem]{Definition}
\newtheorem{algorithm}[theorem]{Algorithm}
\newtheorem{example}[theorem]{Example}
\theoremstyle{remark}
\newtheorem{remark}[theorem]{Remark}
\newtheorem{warning}[theorem]{Warning}

\newtcolorbox{summarybox}[1][]{
  enhanced,breakable,colback=AccentLight,colframe=Accent,
  boxrule=0.7pt,arc=1.5mm,left=2mm,right=2mm,top=1.5mm,bottom=1.5mm,
  title=#1,fonttitle=\bfseries
}
\newtcolorbox{warningbox}[1][]{
  enhanced,breakable,colback=yellow!5,colframe=orange!65!black,
  boxrule=0.7pt,arc=1.5mm,left=2mm,right=2mm,top=1.5mm,bottom=1.5mm,
  title=#1,fonttitle=\bfseries
}

\lstdefinestyle{pythonstyle}{
  language=Python,
  basicstyle=\ttfamily\small,
  keywordstyle=\color{CodeKeyword}\bfseries,
  commentstyle=\color{CodeComment},
  stringstyle=\color{Accent},
  backgroundcolor=\color{CodeBg},
  frame=single,
  rulecolor=\color{RuleGray},
  breaklines=true,
  showstringspaces=false,
  columns=fullflexible,
  keepspaces=true,
  tabsize=4
}

\newcommand{\e}{\mathrm e}
\newcommand{\ii}{\mathrm i}
\newcommand{\dd}{\,\mathrm d}
\newcommand{\W}{W}
\newcommand{\BigO}{\mathcal O}
\newcommand{\LittleO}{o}
\newcommand{\R}{\mathbb R}
\newcommand{\C}{\mathbb C}
\newcommand{\N}{\mathbb N}
\newcommand{\Q}{\mathbb Q}
\newcommand{\K}{\mathbb K}
\newcommand{\calS}{\mathcal S}
\newcommand{\calE}{\mathcal E}
\newcommand{\calR}{\mathcal R}
\newcommand{\calF}{\mathcal F}
\newcommand{\calK}{\mathcal K}
\newcommand{\Gaminv}{\Gamma_{+}^{-1}}
\newcommand{\Barinv}{G_{+}^{-1}}
\newcommand{\Log}{\operatorname{Log}}
\DeclareMathOperator{\val}{val}
\DeclareMathOperator{\supp}{supp}
\DeclareMathOperator{\coeff}{coeff}

\title{\textbf{Asymptotic Transseries for the Inverses of the Gamma and Barnes \(G\)-Functions}\\[0.4em]
\Large Lambert--\(W\) normal forms, all-orders reversion, and residual certification}
\author{Research draft prepared for the ProveIt series-and-transseries program}
\date{3 September 2026}

\begin{document}
\maketitle

\begin{abstract}
We derive all-orders asymptotic transseries for the large positive real inverses of Euler's Gamma function and the Barnes \(G\)-function.  The decisive step is to take logarithms, make a symmetry-adapted shift of the inverse variable, and invert the complete dominant monomial
\[
  A x^{p}(\log x-c)
\]
exactly by the Lambert \(W\)-function.  The remaining inversion is then a near-identity problem over a two-level transseries scale: a slow coefficient field generated by \(w=W(\cdot)\) and \((w+1)^{-1}\), together with fast sectors in powers of the Lambert scale \(X^{-1}\).  We prove a triangular coefficient theorem for locally finite Hahn supports, give a valuation-doubling Newton algorithm, and separate formal coefficient cancellation from the analytic residual-to-error argument needed to justify an actual inverse expansion.

For the principal inverse Gamma branch, put
\[
 M=\log\!\left(\frac{y}{\sqrt{2\pi}}\right),\qquad
 w=W_{0}(M/\e),\qquad T=\frac{M}{w}=\e^{w+1},\qquad s=w+1.
\]
Then \(\Gaminv(y)=\tfrac12+T+\sum_{n\ge1}C_n(s)T^{1-2n}\); the first four nonzero correction sectors are displayed explicitly and a recurrence computes every later coefficient.  The half-shift \(x\mapsto x-\tfrac12\) forces the striking parity law: only odd absolute inverse powers \(T^{-1},T^{-3},\ldots\) occur.

For the eventual positive inverse branch of the Barnes function, put
\[
 Y=\log y-\zeta'(-1),\qquad
 w=W_{0}(4Y\e^{-3}),\qquad T=\sqrt{\frac{4Y}{w}}=\e^{(w+3)/2},\qquad h=w+1.
\]
Then \(\Barinv(y)=1+T+\sum_{n\ge0}D_n(h)T^{-n}\).  We give \(D_0,\ldots,D_5\) explicitly.  The linear term \(\tfrac12\log(2\pi)x\) in the Barnes Stirling expansion creates the leading displacement \(-\log(2\pi)/(w+1)\), while the logarithmic term and Bernoulli tail generate the subsequent sectors.  Numerical substitution into the defining equations confirms the predicted order-by-order decay.
\end{abstract}

\begin{summarybox}[Main formulas at a glance]
\small
For fixed truncation order and \(y\to+\infty\),
\begin{align*}
\Gaminv(y)
 &=\frac12+T+\frac{1}{24sT}
 -\frac{14s^{2}+10s+5}{5760s^{3}T^{3}}\\
 &\quad+\frac{2232s^{4}+1176s^{3}+714s^{2}+350s+105}
 {2903040s^{5}T^{5}}\\
 &\quad-\frac{822960s^{6}+292536s^{5}+136956s^{4}+68040s^{3}
 +32970s^{2}+12250s+2625}
 {1393459200s^{7}T^{7}}
 +\BigO\!\left(\frac{1}{sT^{9}}\right),
\end{align*}
where \(M,w,T,s\) are as in the abstract.  For Barnes \(G\), with
\(c=\log(2\pi)\),
\[
 \Barinv(y)=1+T-\frac{c}{h}+\frac{D_1(h)}{T}
 +\frac{D_2(h)}{T^2}+\frac{D_3(h)}{T^3}
 +\BigO(T^{-4}),
\]
where
\begin{align*}
D_1(h)&=\frac1{12}+\frac1{6h}+c^2\!\left(\frac1{2h^2}-\frac1{h^3}\right),\\
D_2(h)&=\frac{c}{3h^3}+c^3\!\left(\frac{4}{3h^4}-\frac{2}{h^5}\right),\\
D_3(h)&=-\frac1{288}+\frac1{720h}-\frac1{72h^2}-\frac1{36h^3}\\
&\quad+c^2\!\left(-\frac1{24h^2}-\frac1{6h^3}-\frac1{6h^4}+\frac1{h^5}\right)\\
&\quad+c^4\!\left(-\frac1{8h^4}-\frac1{6h^5}+\frac{25}{6h^6}-\frac5{h^7}\right).
\end{align*}
\end{summarybox}

\tableofcontents
\clearpage

\section{Introduction and scope}

\subsection{Which inverse is meant?}

The Gamma function is not one-to-one on the whole real line and its complex inverse is infinitely multivalued.  Throughout the main development,
\[
  \Gaminv(y)
\]
denotes the inverse of the restriction of \(\Gamma(x)\) to its eventual increasing positive-real branch.  Equivalently, for all sufficiently large \(y>0\), \(\Gaminv(y)\) is the unique large positive solution of \(\Gamma(x)=y\).  This is the principal real inverse in the standard terminology.  It must not be confused with the inverse-Gamma probability distribution.  Real and complex inverse branches, their analytic structure, and numerical computation are treated in \cite{Uchiyama,Pedersen,FolitseJeffreyCorless,JeffreyWatt}; related Lambert--\(W\) approximations are discussed in \cite{BorweinCorless}.

Similarly, \(\Barinv(y)\) denotes the unique large positive solution of
\(G(x)=y\), where \(G\) is the Barnes \(G\)-function normalized by
\(G(1)=1\) and \(G(z+1)=\Gamma(z)G(z)\).  We deliberately say
``eventual positive inverse branch'': no global monotonicity assertion is needed.  The derivative of the large-\(x\) logarithmic asymptotic is positive, so existence and uniqueness on a sufficiently far positive ray follow directly.  The normalization goes back to Barnes \cite{Barnes1900}; standard formulas are collected in \cite{DLMF}.

The phrase ``large argument'' means \(y\to+\infty\) in the inverse equation.  Since the functions themselves grow superexponentially in the appropriate sense, the inverse variable grows much more slowly: roughly \(\log y/\log\log y\) for Gamma and \(\sqrt{\log y/\log\log y}\) for Barnes \(G\).

\subsection{Why ordinary power series are the wrong language}

After taking logarithms, Gamma inversion begins with
\[
  t\log t-t=M,
\]
while Barnes inversion begins with
\[
  \frac12t^2\log t-\frac34t^2=Y.
\]
Neither inverse has an ordinary Laurent expansion in \(M^{-1}\) or \(Y^{-1}\).  The leading answers contain nested logarithms, and the cleanest exact coordinates are Lambert--\(W\) coordinates.  Once the dominant balance has been inverted exactly, the correction terms are powers of an exponentially large scale \(T\), whose coefficients are rational functions of the slow variable \(w=W(\cdot)\).  Expanding \(w\) itself creates a second logarithmic level \(\log\log\), and expanding the coefficients creates powers of \(1/\log\).  Thus the natural object is a \emph{nested transseries}, not a single series.

This point agrees with the general philosophy developed in the ProveIt polynomial--logarithmic transseries draft: composition and reversal must preserve both the power scale and all logarithmic generators, and the full leading monomial---not merely its power part---must be normalized before near-identity reversion is attempted \cite{ProveItPLT,SalvyShackell,vanderHoeven}.

\subsection{Contributions}

The article makes five concrete contributions.

\begin{enumerate}[label=\textup{(\roman*)}]
\item It converts the Gamma and Barnes Stirling expansions into symmetry-adapted normal forms in which the dominant inverse is exact in \(W_0\).
\item It derives explicit all-orders recurrences and gives several high-order correction coefficients in compact rational form.
\item It proves a general inversion theorem for
\(A x^p(\log x-c)\) plus a Hahn-supported perturbation, including support propagation and coefficient-ring closure.
\item It supplies two independent computational mechanisms: triangular coefficient extraction and formal Newton iteration with valuation doubling.
\item It gives the analytic step often omitted in purely formal manipulations: a residual divided by a derivative lower bound controls the error of the actual inverse.
\end{enumerate}

\subsection{A map of the argument}

\Cref{sec:inputs} records and simplifies the asymptotic inputs.  The exact Lambert normal form is proved in \cref{sec:dominant}.  The general formal and analytic inversion theory occupies \cref{sec:general}.  The Gamma and Barnes calculations appear in \cref{sec:gamma,sec:barnes}.  Log--log re-expansions, complex branches, Stokes sectors, numerical tests, and reproducible symbolic code follow.

\section{Asymptotic input data}\label{sec:inputs}

We use Bernoulli numbers and polynomials in the conventions
\[
 \frac{z}{\e^z-1}=\sum_{n\ge0}B_n\frac{z^n}{n!},
 \qquad
 \frac{z\e^{xz}}{\e^z-1}=\sum_{n\ge0}B_n(x)\frac{z^n}{n!}.
\]
Thus \(B_1=-\tfrac12\), \(B_2=\tfrac16\), \(B_4=-\tfrac1{30}\), and
\[
 B_n\!\left(\frac12\right)=\bigl(2^{1-n}-1\bigr)B_n.
\]

\subsection{The shifted Gamma expansion}

The sectorial Stirling expansions used here are standard; see \cite{DLMF,Olver}.  The usual Stirling expansion contains the leading block
\((x-\tfrac12)\log x-x\).  For inversion, a half-unit translation is substantially better.

\begin{lemma}[Symmetry-adapted Gamma normal form]\label{lem:gamma-stirling}
Let \(t=x-\tfrac12\).  On every closed subsector of
\(|\arg t|<\pi\),
\begin{equation}\label{eq:gamma-shifted}
 \log\Gamma\!\left(t+\frac12\right)
 \sim t\log t-t+\frac12\log(2\pi)
 +\sum_{n\ge1}\frac{a_n}{t^{2n-1}},
\end{equation}
where
\begin{equation}\label{eq:an}
 a_n=\frac{B_{2n}(1/2)}{2n(2n-1)}
 =\frac{(2^{1-2n}-1)B_{2n}}{2n(2n-1)}.
\end{equation}
In particular,
\[
 a_1=-\frac1{24},\quad
 a_2=\frac7{2880},\quad
 a_3=-\frac{31}{40320},\quad
 a_4=\frac{127}{215040},\quad
 a_5=-\frac{511}{608256}.
\]
Only odd inverse powers of \(t\) occur.
\end{lemma}

\begin{proof}
The shifted Stirling expansion is
\[
 \log\Gamma(t+h)
 \sim \left(t+h-\frac12\right)\log t-t+\frac12\log(2\pi)
 +\sum_{k\ge2}\frac{(-1)^kB_k(h)}{k(k-1)t^{k-1}}.
\]
Set \(h=\tfrac12\).  The coefficient of \(\log t\) becomes exactly \(t\), and
\(B_{2m+1}(1/2)=0\) for \(m\ge1\).  Relabel \(k=2n\) to obtain
\eqref{eq:gamma-shifted}--\eqref{eq:an}.
\end{proof}

\begin{remark}[Why the half-shift matters]
Had one expanded directly in \(x^{-1}\), both parities would appear and the inverse recurrence would be needlessly dense.  The translation by \(\tfrac12\) exposes the reflection-centered structure of Stirling's series and ultimately forces the inverse correction powers \(T^{-1},T^{-3},T^{-5},\ldots\).
\end{remark}

\subsection{The Barnes \texorpdfstring{\(G\)}{G}-expansion in inversion form}

The standard Barnes formula can be combined with Stirling's expansion for
\(\log\Gamma(t+1)\).  The resulting expression is much better suited to inversion; compare \cite{DLMF,FerreiraLopez,NemesBarnes}.

\begin{lemma}[Simplified Barnes normal form]\label{lem:barnes-stirling}
As \(t\to\infty\) in every closed subsector of \(|\arg t|<\pi\),
\begin{equation}\label{eq:barnes-simplified}
\begin{split}
 \log G(t+1)
 &\sim \frac12t^2\log t-\frac34t^2
 +\frac{c}{2}t-\frac1{12}\log t+\zeta'(-1)
 +\sum_{k\ge1}\frac{b_k}{t^{2k}},
\end{split}
\end{equation}
where
\begin{equation}\label{eq:bk}
 c=\log(2\pi),\qquad
 b_k=\frac{B_{2k+2}}{4k(k+1)}.
\end{equation}
Thus
\[
 b_1=-\frac1{240},\qquad
 b_2=\frac1{1008},\qquad
 b_3=-\frac1{1440},\qquad
 b_4=\frac1{1056}.
\]
\end{lemma}

\begin{proof}
Start from
\[
\begin{split}
 \log G(t+1)
 &\sim \frac14t^2+t\log\Gamma(t+1)
 -\left(\frac12t(t+1)+\frac1{12}\right)\log t-\log A\\
 &\quad+\sum_{k\ge1}
 \frac{B_{2k+2}}{2k(2k+1)(2k+2)t^{2k}},
\end{split}
\]
where \(A\) is the Glaisher--Kinkelin constant.  Insert
\[
 \log\Gamma(t+1)
 \sim \left(t+\frac12\right)\log t-t+\frac12\log(2\pi)
 +\sum_{j\ge1}\frac{B_{2j}}{2j(2j-1)t^{2j-1}}.
\]
The quadratic and logarithmic terms simplify immediately.  At inverse power
\(t^{-2k}\), the Gamma contribution with \(j=k+1\) and the explicit Barnes contribution add to
\[
 \frac{B_{2k+2}}{(2k+2)(2k+1)}
 +\frac{B_{2k+2}}{2k(2k+1)(2k+2)}
 =\frac{B_{2k+2}}{4k(k+1)}.
\]
The Gamma term with \(j=1\) contributes \(B_2/2=1/12\), while
\(\log A=1/12-\zeta'(-1)\); the constants therefore combine to \(\zeta'(-1)\).
\end{proof}

\begin{warningbox}[Poincar\'e series versus exponentially improved expansions]
The series in \cref{eq:gamma-shifted,eq:barnes-simplified} are Poincar\'e asymptotic series.  They are generally divergent.  Fixed-order inverse expansions are rigorously meaningful after a fixed truncation, but optimal truncation and exponentially small terms require a larger transseries.  For Barnes \(G\), the Stokes lines of the logarithmic expansion occur at \(\arg t=\pm\pi/2\); this issue is revisited in \cref{sec:stokes}.
\end{warningbox}

\section{Exact inversion of the dominant logarithmic monomial}\label{sec:dominant}

The same exact calculation underlies both special functions.  We use the standard branch conventions for Lambert's function \cite{CorlessW,DLMF}.

\begin{theorem}[Lambert normal form]\label{thm:lambert-normal}
Let \(A>0\), \(p>0\), and \(c\in\R\), and define
\[
 \Phi_0(x)=A x^p(\log x-c),\qquad x>\e^c.
\]
For \(Y>0\), put
\begin{equation}\label{eq:general-scales}
 U=\frac{pY}{A},\qquad
 w=W_0\!\left(U\e^{-pc}\right),\qquad
 X=\left(\frac{U}{w}\right)^{1/p}.
\end{equation}
Then
\begin{equation}\label{eq:dominant-exact}
 \Phi_0(X)=Y,
 \qquad
 \log X=c+\frac{w}{p},
 \qquad
 \Phi_0'(X)=A X^{p-1}(w+1).
\end{equation}
The same identities hold complex-sectorially with a consistent choice of
\(W_k\), logarithm, and \(p\)-th root.
\end{theorem}

\begin{proof}
By definition,
\[
 w\e^w=U\e^{-pc}.
\]
Since \(X^p=U/w\), taking logarithms gives
\[
 p\log X=\log U-\log w=w+pc,
\]
which proves \(\log X=c+w/p\).  Therefore
\[
 \Phi_0(X)=A\frac{U}{w}\frac{w}{p}=Y.
\]
Finally,
\[
 \Phi_0'(x)=A x^{p-1}\bigl(p(\log x-c)+1\bigr),
\]
and evaluation at \(X\) gives the last identity.
\end{proof}

\begin{remark}[The universal divisor \(w+1\)]
Every inverse correction is obtained by dividing a known residual coefficient by \(w+1\).  This factor is not an algebraic accident: it is the normalized derivative of the full leading monomial.  Its zero at \(w=-1\) is the Lambert branch point, where ordinary inverse conditioning fails.
\end{remark}

\subsection{The normalized displacement equation}

Write an exact or asymptotic equation in the form
\begin{equation}\label{eq:perturbed-general}
 \Phi(x)=A x^p(\log x-c)+A x^p R(x)=Y,
\end{equation}
and set
\[
 x=X(1+u),
\]
where \(X\) is the exact dominant inverse from \cref{thm:lambert-normal}.  Dividing by \(AX^p\) gives
\begin{equation}\label{eq:universal-normalized}
 H_{p,w}(u)+\mathscr R_X(u)=0,
\end{equation}
where
\begin{align}
 H_{p,w}(u)
 &= (1+u)^p\left(\frac{w}{p}+\log(1+u)\right)-\frac{w}{p},
 \label{eq:Hpw}\\
 \mathscr R_X(u)
 &= (1+u)^p R\bigl(X(1+u)\bigr).
 \label{eq:RX}
\end{align}
The key local expansion is
\begin{equation}\label{eq:Hlinear}
 H_{p,w}(u)=(w+1)u+\BigO(u^2).
\end{equation}
For \(p=1\), the nonlinear part is particularly simple:
\begin{equation}\label{eq:H-p1}
 H_{1,w}(u)=(w+1)u+
 \sum_{m\ge2}\frac{(-1)^m}{m(m-1)}u^m.
\end{equation}

\section{A general inversion theory for Lambert-shaped transseries}\label{sec:general}

The valued-series viewpoint is compatible with the inverse-multiseries methods of Salvy and Shackell and with the broader transseries framework of van der Hoeven \cite{SalvyShackell,vanderHoeven}.

\subsection{Coefficient fields and admissible supports}

Let \(w\) be a slow parameter tending to infinity in a fixed sector, and let
\(X=\exp(c+w/p)\) be the fast Lambert scale.  We use a coefficient field
\(\calK\) containing rational functions of \(w\) and \((w+1)^{-1}\); in applications it also contains constants such as \(\log(2\pi)\), Bernoulli numbers, and \(\zeta'(-1)\).

Let \(\Lambda\subset(0,\infty)\) be a set of perturbation exponents, and let
\[
 \calS=\left\{\lambda_1+\cdots+\lambda_m:
 m\ge0,\ \lambda_j\in\Lambda\right\}
\]
be its additive semigroup.  We assume \(\calS\) is \emph{locally finite}: every bounded interval contains only finitely many elements of \(\calS\).  Finite sets of positive rational exponents, and more generally well-ordered Hahn supports with the finite-decomposition property, satisfy the requirement.

A formal fast-scale transseries is
\[
 U(X)=\sum_{\mu\in\calS}U_\mu(w)X^{-\mu},
 \qquad U_\mu\in\calK,
\]
with valuation
\[
 \val(U)=\min\{\mu:U_\mu\ne0\}.
\]
Local finiteness makes addition, Cauchy multiplication, and coefficient extraction finite at every prescribed exponent.

\subsection{Triangular inversion}

Consider a perturbation
\begin{equation}\label{eq:R-hahn}
 R(x)=\sum_{\lambda\in\Lambda}
 r_\lambda(\log x)x^{-\lambda},
\end{equation}
where each \(r_\lambda\) belongs to a coefficient algebra stable under differentiation and translation by a formal logarithm.  Substitution into \cref{eq:RX} gives
\begin{equation}\label{eq:RX-expanded}
 \mathscr R_X(u)=
 \sum_{\lambda\in\Lambda}X^{-\lambda}
 (1+u)^{p-\lambda}
 r_\lambda\!\left(c+\frac{w}{p}+\log(1+u)\right).
\end{equation}

\begin{theorem}[Triangular Lambert-scale inversion]\label{thm:triangular-general}
Assume \(w+1\) is invertible in \(\calK\), \(\calS\) is locally finite, and the substitutions in \cref{eq:RX-expanded} are formally defined.  Then the normalized equation
\[
 H_{p,w}(u)+\mathscr R_X(u)=0
\]
has a unique solution of positive valuation,
\begin{equation}\label{eq:u-hahn}
 u=\sum_{\mu\in\calS\setminus\{0\}}U_\mu(w)X^{-\mu}.
\end{equation}
If \(\mu_1=\min\Lambda\), then
\begin{equation}\label{eq:first-general-coeff}
 U_{\mu_1}(w)
 =-\frac{r_{\mu_1}(c+w/p)}{w+1}.
\end{equation}
More generally, if \(u_{<\mu}\) denotes the truncation to exponents below \(\mu\), then
\begin{equation}\label{eq:general-tri-rec}
 U_\mu(w)
 =-\frac{1}{w+1}
 [X^{-\mu}]
 \left(
 H_{p,w}(u_{<\mu})-(w+1)u_{<\mu}
 +\mathscr R_X(u_{<\mu})
 \right).
\end{equation}
\end{theorem}

\begin{proof}
By \cref{eq:Hlinear}, the coefficient of the as-yet unknown block
\(U_\mu X^{-\mu}\) in the left-hand side is exactly
\((w+1)U_\mu X^{-\mu}\).  Every other contribution at exponent \(\mu\) is a finite sum of products of blocks whose exponents are strictly smaller than \(\mu\), together with a direct perturbation block.  Local finiteness guarantees that only finitely many decompositions of \(\mu\) occur.  Hence \cref{eq:general-tri-rec} determines \(U_\mu\) uniquely, inductively through the ordered support.  At the smallest exponent, no nonlinear product can contribute, yielding \cref{eq:first-general-coeff}.
\end{proof}

\begin{corollary}[Coefficient-ring closure]\label{cor:coefficient-ring}
Suppose \(p\in\Q_{>0}\), \(c\in K\), every exponent in \(\Lambda\) is rational, and every
\(r_\lambda(L)\) is a polynomial over a field \(K\) of characteristic zero.  Then every inverse coefficient belongs to
\[
 K\bigl[w,(w+1)^{-1}\bigr].
\]
If rational functions of \(L\) are allowed, the coefficient field must also invert their translated denominators.
\end{corollary}

\begin{proof}
Binomial expansion, Taylor expansion in \(\log(1+u)\), and multiplication use only rational operations and derivatives of the polynomial coefficients.  The sole new division in each triangular step is by \(w+1\).
\end{proof}

\subsection{A coefficient formula on a single lattice}

Many applications reduce to integer powers of one small parameter \(\varepsilon\).  Suppose
\begin{equation}\label{eq:single-lattice-E}
 \calE(u,\varepsilon)
 =a u+\sum_{j\ge2}h_j u^j
 +\sum_{k\ge1}\varepsilon^k\sum_{j\ge0}r_{k,j}u^j=0,
 \qquad a\ne0,
\end{equation}
with \(u=\sum_{n\ge1}u_n\varepsilon^n\).  Then
\begin{equation}\label{eq:bell-recurrence}
\begin{split}
 u_n=-\frac1a\Bigg(&
 \sum_{j=2}^{n}h_j
 \sum_{\substack{n_1+\cdots+n_j=n\\n_i\ge1}}
 u_{n_1}\cdots u_{n_j}\\
 &+\sum_{k=1}^{n}\sum_{j=0}^{n-k}r_{k,j}
 \sum_{\substack{n_1+\cdots+n_j=n-k\\n_i\ge1}}
 u_{n_1}\cdots u_{n_j}
 \Bigg),
\end{split}
\end{equation}
where the inner sum for \(j=0\) is \(1\) when \(n=k\) and \(0\) otherwise.  The convolution sums can be expressed through ordinary partial Bell polynomials.  Formula \cref{eq:bell-recurrence} is the direct, implementation-friendly form.

When the equation has the special shape
\[
 H(u)+\varepsilon R(u)=0,
 \qquad H(0)=0,\quad H'(0)\ne0,
\]
one may write
\[
 u=\varepsilon\phi(u),
 \qquad
 \phi(u)=-\frac{uR(u)}{H(u)},
\]
and Lagrange--B\"urmann gives
\begin{equation}\label{eq:lagrange-special}
 [\varepsilon^n]u(\varepsilon)
 =\frac1n[z^{n-1}]\phi(z)^n.
\end{equation}
The triangular recurrence is more flexible because it handles several fast exponents and explicit \(\varepsilon\)-dependence without diagonal extraction.

\subsection{Formal Newton iteration}

\begin{theorem}[Valuation doubling]\label{thm:newton-doubling}
Let \(\calE(u)=H_{p,w}(u)+\mathscr R_X(u)\), and suppose
\(\calE'(u)\) has leading coefficient \(w+1\), hence is a unit of valuation zero.  Define
\begin{equation}\label{eq:newton-formal}
 u^{[m+1]}=u^{[m]}-\frac{\calE(u^{[m]})}{\calE'(u^{[m]})}.
\end{equation}
If \(\val(\calE(u^{[m]}))\ge\rho\), then
\[
 \val(\calE(u^{[m+1]}))\ge2\rho.
\]
Thus a precision-truncated implementation approximately doubles the number of correct fast-scale blocks at each Newton step.
\end{theorem}

\begin{proof}
Put \(\Delta=-\calE/\calE'\).  Formal Taylor expansion gives
\[
 \calE(u+\Delta)=\calE(u)+\calE'(u)\Delta
 +\sum_{j\ge2}\frac{\calE^{(j)}(u)}{j!}\Delta^j.
\]
The first two terms cancel.  Since \(\calE'\) is a valuation-zero unit,
\(\val(\Delta)=\val(\calE(u))\ge\rho\), and every remaining term has valuation at least \(2\rho\).
\end{proof}

\begin{remark}[Triangular recurrence versus Newton]
The recurrence exposes exact coefficient formulas and is ideal for a paper.  Newton iteration is usually faster in computer algebra at high order.  The two methods also provide an independent cross-check: their truncated residuals must agree through the requested valuation.
\end{remark}

\subsection{From a formal residual to an error of the actual inverse}

Formal cancellation alone does not prove an asymptotic statement about an analytic function.  The missing step is elementary but essential.

\begin{theorem}[Residual-to-error transfer]\label{thm:residual-error}
Let \(F\) be continuously differentiable and strictly monotone on an interval
\(I\).  Let \(x_*\in I\) satisfy \(F(x_*)=Y\), and let \(x_N\in I\) be an approximation.  If
\[
 m_N=\inf_{\xi\text{ between }x_N\text{ and }x_*}|F'(\xi)|>0,
\]
then
\begin{equation}\label{eq:residual-error}
 |x_N-x_*|\le \frac{|F(x_N)-Y|}{m_N}.
\end{equation}
For a Lambert-normalized family with \(x_N\sim X\),
\[
 F'(x_N)\sim A X^{p-1}(w+1).
\]
Consequently, a normalized residual
\(\BigO(X^{-\rho})\) in \cref{eq:universal-normalized} produces an inverse error
\begin{equation}\label{eq:order-transfer}
 x_N-x_*=\BigO\!\left(\frac{X^{1-\rho}}{w+1}\right)
\end{equation}
for fixed truncation order, provided the asymptotic input remainder has the corresponding analytic bound.
\end{theorem}

\begin{proof}
The mean value theorem gives
\[
 F(x_N)-F(x_*)=F'(\xi)(x_N-x_*)
\]
for some intermediate \(\xi\), proving \cref{eq:residual-error}.  The derivative asymptotic follows from \cref{eq:dominant-exact} and the fact that \(x_N/X\to1\).  Multiplying a normalized residual \(\BigO(X^{-\rho})\) by \(AX^p\), then dividing by the derivative scale \(AX^{p-1}(w+1)\), yields \cref{eq:order-transfer}.
\end{proof}

\begin{warningbox}[A formal order is not an analytic bound]
To use \cref{thm:residual-error}, one must know that the truncated Stirling expansion has a remainder of the claimed size in the chosen sector.  This is true for each fixed truncation on positive real rays and on closed sectors away from the cut, but it is a separate theorem about the special function.  Differentiating an arbitrary asymptotic series term-by-term is likewise not automatically valid without a suitable realization theorem or derivative estimate.
\end{warningbox}

\subsection{Why this is a genuine two-level transseries}

Because
\[
 X=\exp\!\left(c+\frac{w}{p}\right),
\]
every fast monomial \(X^{-\lambda}=\exp(-\lambda c-\lambda w/p)\) is smaller than every algebraic power \(w^{-N}\).  Hence the inverse has the hierarchical form
\begin{equation}\label{eq:two-level}
 x=X\left(1+
 \sum_{\mu>0}U_\mu(w)\e^{-\mu(c+w/p)}\right),
\end{equation}
where each \(U_\mu(w)\) can itself be expanded as a Laurent series in \(w^{-1}\), and \(w\) can be expanded in nested logarithms of the original argument.  The three nested levels are therefore:
\[
 \begin{aligned}
 \text{original argument}&\Longrightarrow \log,\log\log\text{ sector}
 \Longrightarrow w^{-1}\text{ coefficients}\\
 &\Longrightarrow \e^{-\mu w/p}\text{ fast sectors}.
 \end{aligned}
\]
Keeping \(w\) exact is both shorter and numerically more stable; expanding it is useful when an answer solely in elementary logarithms is required.

\section{The inverse Gamma transseries}\label{sec:gamma}

\subsection{Lambert coordinates}

Let \(y\to+\infty\), and define
\begin{equation}\label{eq:gamma-scales}
 M=\log\!\left(\frac{y}{\sqrt{2\pi}}\right),\qquad
 w=W_0(M/\e),\qquad
 T=\frac{M}{w}=\e^{w+1},\qquad
 s=w+1=\log T.
\end{equation}
The identities follow from \(w\e^w=M/\e\).  If \(t=x-\tfrac12\), then the dominant equation in \cref{eq:gamma-shifted} is
\[
 t\log t-t=M.
\]
It is solved exactly by \(t=T\).

Write
\begin{equation}\label{eq:gamma-u}
 t=T(1+u),\qquad \varepsilon=T^{-2}.
\end{equation}
After subtracting the dominant equation and dividing by \(T\), the complete formal inversion equation is
\begin{equation}\label{eq:gamma-rec-equation}
 0=su+
 \sum_{m\ge2}\frac{(-1)^m}{m(m-1)}u^m
 +\sum_{n\ge1}a_n\varepsilon^n(1+u)^{-(2n-1)}.
\end{equation}
This equation makes both the parity and the triangular structure transparent.

\begin{theorem}[All-orders inverse-Gamma transseries]\label{thm:gamma-all-orders}
For every fixed \(N\ge0\), as \(y\to+\infty\),
\begin{equation}\label{eq:gamma-all-orders}
 \Gaminv(y)=\frac12+T+
 \sum_{n=1}^{N}C_n(s)T^{1-2n}
 +\mathcal O_N\!\left(\frac{T^{-2N-1}}{s}\right),
\end{equation}
where \(M,w,T,s\) are given by \cref{eq:gamma-scales}.  The coefficients are the unique solution of
\begin{equation}\label{eq:gamma-C-recurrence}
 u(\varepsilon)=\sum_{n\ge1}C_n(s)\varepsilon^n
\end{equation}
substituted into \cref{eq:gamma-rec-equation}.  Equivalently,
\begin{equation}\label{eq:gamma-C-triangular}
 C_n(s)=-\frac1s[\varepsilon^n]
 \left\{
 \sum_{m\ge2}\frac{(-1)^m}{m(m-1)}u_{<n}^{m}
 +\sum_{j\ge1}a_j\varepsilon^j(1+u_{<n})^{-(2j-1)}
 \right\},
\end{equation}
where \(u_{<n}=\sum_{j=1}^{n-1}C_j(s)\varepsilon^j\).  Moreover,
\begin{equation}\label{eq:gamma-C-ring}
 C_n(s)\in\Q[s^{-1}],\qquad C_n(s)=\BigO_n(s^{-1})
 \quad(s\to+\infty).
\end{equation}
Only the absolute powers \(T^{-1},T^{-3},T^{-5},\ldots\) occur.
\end{theorem}

\begin{proof}
\Cref{eq:gamma-rec-equation} has linear coefficient \(s=w+1\), which is positive on the principal real branch.  The triangular theorem therefore gives a unique formal series in \(\varepsilon=T^{-2}\).  Since the input coefficients \(a_n\) are rational and the only division is by \(s\), induction gives \(C_n\in\Q[s^{-1}]\).  The right-hand side that determines \(sC_n\) remains bounded as \(s\to\infty\): every earlier \(C_j\) is \(O(s^{-1})\), and the direct term \(a_n\) is constant.  This proves \(C_n=O(s^{-1})\).

Because \(t=T(1+u)\), the term \(C_n\varepsilon^n\) in \(u\) becomes the absolute correction \(C_nT^{1-2n}\).  No other exponent can occur.  Finally, the fixed-order remainder in the shifted Stirling expansion, together with \cref{thm:residual-error}, transfers the first omitted normalized order \(T^{-2N-2}\) to the inverse error \(O(T^{-2N-1}/s)\).
\end{proof}

\subsection{Explicit coefficients}

Straight coefficient extraction from \cref{eq:gamma-rec-equation} gives
\begin{align}
 C_1(s)&=\frac{1}{24s},\label{eq:C1}\\
 C_2(s)&=-\frac{14s^2+10s+5}{5760s^3},\label{eq:C2}\\
 C_3(s)&=\frac{2232s^4+1176s^3+714s^2+350s+105}
 {2903040s^5},\label{eq:C3}\\
 C_4(s)&=-\frac{
 822960s^6+292536s^5+136956s^4+68040s^3
 +32970s^2+12250s+2625}
 {1393459200s^7},\label{eq:C4}\\
 C_5(s)&=\frac{
 309052800s^8+77920128s^7+26024592s^6+10019592s^5
 +4516028s^4}
 {367873228800s^9}\notag\\
 &\quad+\frac{
 2023560s^3+812350s^2+242550s+40425}
 {367873228800s^9}.
 \label{eq:C5}
\end{align}
The split line in \cref{eq:C5} is typographical; the two numerators have the same denominator.

Thus a five-sector version is
\begin{equation}\label{eq:gamma-five-sector}
 \Gaminv(y)=\frac12+T+\frac{C_1(s)}{T}
 +\frac{C_2(s)}{T^3}+\frac{C_3(s)}{T^5}
 +\frac{C_4(s)}{T^7}+\frac{C_5(s)}{T^9}
 +\BigO\!\left(\frac{T^{-11}}{s}\right).
\end{equation}

\subsection{Deriving the first coefficients by hand}

It is useful to see the triangular mechanism explicitly.  Put
\[
 u=C_1\varepsilon+C_2\varepsilon^2+C_3\varepsilon^3+\cdots.
\]
At order \(\varepsilon\), \cref{eq:gamma-rec-equation} gives
\[
 sC_1+a_1=0,
\]
so \(C_1=-a_1/s=1/(24s)\).  At order \(\varepsilon^2\),
\[
 sC_2+\frac12C_1^2-a_1C_1+a_2=0.
\]
Substituting \(a_1=-1/24\), \(a_2=7/2880\), and \(C_1=1/(24s)\) yields \cref{eq:C2}.  At order \(\varepsilon^3\),
\[
\begin{split}
0={}&sC_3+C_1C_2-\frac16C_1^3
 +a_1(C_1^2-C_2)-3a_2C_1+a_3,
\end{split}
\]
which simplifies to \cref{eq:C3}.  All later orders are of the same triangular type; their apparent size comes from convolution, not from any conceptual complication.

\subsection{Elementary log--log expansion of the Lambert scale}\label{sec:gamma-loglog}

Some applications require eliminating \(W_0\) in favor of elementary logarithms.  Put
\begin{equation}\label{eq:gamma-log-vars}
 A=\log M,
 \qquad B=\log A,
 \qquad D=B+1.
\end{equation}
The exact relation \(T(\log T-1)=M\) admits the expansion
\begin{equation}\label{eq:T-gamma-loglog}
\begin{split}
 T=\frac{M}{A}\Bigg[&1+\frac{D}{A}
 +\frac{D(D-1)}{A^2}
 +\frac{D(D-2)(2D-1)}{2A^3}\\
 &+\frac{D(6D^3-26D^2+27D-6)}{6A^4}\\
 &+\frac{D(12D^4-77D^3+142D^2-84D+12)}{12A^5}
 +\BigO\!\left(\frac{B^6}{A^6}\right)\Bigg].
\end{split}
\end{equation}

\begin{proof}[Derivation]
Write \(T=(M/A)F\), where \(F=1+\sum_{n\ge1}f_nA^{-n}\).  Since
\(\log T=A-B+\log F\), the defining equation is equivalent to
\begin{equation}\label{eq:Fgamma}
 F\left(1-\frac{D}{A}+\frac{\log F}{A}\right)=1.
\end{equation}
Equating powers of \(A^{-1}\) gives
\[
 f_1=D,
 \quad f_2=D(D-1),
 \quad f_3=\frac12D(D-2)(2D-1),
\]
and then the displayed higher coefficients.  This is an ordinary formal reversal after the correct nested logarithm \(B=\log A\) has been introduced.
\end{proof}

Combining \cref{eq:T-gamma-loglog} with \cref{eq:gamma-all-orders} gives a fully elementary nested-log transseries.  It is important, however, to understand the hierarchy.  Since \(T\asymp M/A\), every correction \(T^{1-2n}\) is beyond all orders of the relative \(A^{-1}\) expansion of \(T\).  No finite extension of \cref{eq:T-gamma-loglog} can reproduce even the first \(1/(24sT)\) sector.

\begin{summarybox}[Gamma hierarchy]
The principal inverse Gamma transseries naturally separates as
\[
 \Gaminv(y)
 =\frac12+
 \underbrace{T(M)}_{\text{nested-log sector}}
 +\underbrace{\frac{C_1(s)}{T}+\frac{C_2(s)}{T^3}+\cdots}_{\text{beyond all relative log orders}},
\]
where each \(C_n(s)\) may itself be expanded in \(s^{-1}\), and \(s=w+1\) may be expanded in \(A^{-1}\) with polynomial coefficients in \(B=\log A\).
\end{summarybox}

\subsection{A direct residual certificate}

Let
\[
 x_N=\frac12+T+\sum_{n=1}^{N}C_n(s)T^{1-2n}.
\]
A symbolic implementation should not merely print \(C_n\); it should verify the certificate
\begin{equation}\label{eq:gamma-residual-cert}
 \log\Gamma(x_N)-\log y
 =\BigO\!\left(T^{-2N-1}\right),
\end{equation}
when the logarithmic Gamma expansion is retained to compatible order.  Since
\[
 \frac{\dd}{\dd x}\log\Gamma(x)=\psi(x)\sim\log x\sim s,
\]
\cref{eq:gamma-residual-cert} converts to
\[
 x_N-\Gaminv(y)=\BigO\!\left(\frac{T^{-2N-1}}{s}\right),
\]
matching \cref{thm:gamma-all-orders}.  Using the logarithmic residual is numerically preferable to evaluating \(\Gamma(x_N)-y\), which can overflow long before the inverse approximation loses accuracy.

\subsection{Complex branch template}\label{sec:gamma-complex}

For complex \(y\), taking logarithms introduces a sheet index.  On a chosen logarithm sheet, define
\[
 M_m=\Log y+2\pi\ii m-\frac12\log(2\pi),
 \qquad m\in\mathbb Z.
\]
For a Lambert branch \(k\in\mathbb Z\), set
\[
 w_{m,k}=W_k(M_m/\e),
 \qquad T_{m,k}=M_m/w_{m,k},
 \qquad s_{m,k}=w_{m,k}+1.
\]
Whenever the resulting \(T_{m,k}\) remains in a sector where the shifted Stirling expansion is valid and all branch choices are continued consistently, the same formal coefficient functions give
\begin{equation}\label{eq:gamma-complex-template}
 x_{m,k}\sim\frac12+T_{m,k}
 +\sum_{n\ge1}C_n(s_{m,k})T_{m,k}^{1-2n}.
\end{equation}
This is a local sectorial template, not a global classification of the Riemann surface of \(\Gamma^{-1}\); for the global computational geometry, see \cite{JeffreyWatt}.  Different pairs \((m,k)\) can describe continuations that meet, and critical points of \(\Gamma\) create additional branch geometry not visible in the leading Lambert model.

\section{The inverse Barnes \texorpdfstring{\(G\)}{G}-transseries}\label{sec:barnes}

\subsection{Lambert coordinates for the quadratic-logarithmic scale}

Let
\begin{equation}\label{eq:barnes-scales}
 Y=\log y-\zeta'(-1),\qquad
 w=W_0(4Y\e^{-3}),\qquad
 T=\sqrt{\frac{4Y}{w}}=\e^{(w+3)/2},\qquad
 h=w+1.
\end{equation}
Then
\[
 \log T=1+\frac h2=\frac{w+3}{2}.
\]
If \(t=x-1\), the dominant Barnes equation is
\begin{equation}\label{eq:barnes-dominant}
 \frac12t^2\log t-\frac34t^2=Y,
\end{equation}
and \(t=T\) solves it exactly.

Set
\begin{equation}\label{eq:barnes-delta}
 t=T+\delta,\qquad \varepsilon=T^{-1},\qquad
 \delta(\varepsilon)=\sum_{n\ge0}D_n(h)\varepsilon^n,
\end{equation}
and abbreviate
\[
 c=\log(2\pi),\qquad r=\frac h2,
 \qquad \sigma=1+\frac h2=\log T,
 \qquad a=\frac c2.
\]
Subtracting \cref{eq:barnes-dominant}, inserting \cref{eq:barnes-simplified}, and dividing by \(T\) gives
\begin{equation}\label{eq:barnes-rec-equation}
\begin{split}
0={}&r\delta+a
 +\varepsilon\left(\frac{\sigma}{2}\delta^2+a\delta-\frac{\sigma}{12}\right)\\
&+\sum_{m\ge3}\frac{(-1)^{m-3}}{m(m-1)(m-2)}
 \delta^m\varepsilon^{m-1}
 -\frac{\varepsilon}{12}\log(1+\varepsilon\delta)\\
&+\sum_{k\ge1}b_k\varepsilon^{2k+1}
 (1+\varepsilon\delta)^{-2k}.
\end{split}
\end{equation}
Unlike the Gamma equation, this recurrence begins with a nonzero constant displacement because the Barnes Stirling expansion contains a linear term \(ct/2\).

\begin{theorem}[All-orders inverse-Barnes transseries]\label{thm:barnes-all-orders}
For every fixed \(N\ge0\), as \(y\to+\infty\),
\begin{equation}\label{eq:barnes-all-orders}
 \Barinv(y)=1+T+\sum_{n=0}^{N}D_n(h)T^{-n}
 +\mathcal O_N(T^{-N-1}),
\end{equation}
where \(Y,w,T,h\) are given by \cref{eq:barnes-scales}.  The coefficients are determined uniquely by \cref{eq:barnes-rec-equation}; explicitly,
\begin{equation}\label{eq:barnes-triangular}
 D_n(h)=-\frac{2}{h}[\varepsilon^n]
 \calE(\delta_{<n},\varepsilon),
\end{equation}
where \(\calE\) denotes the right-hand side of
\cref{eq:barnes-rec-equation} with the term \(r\delta\) included, and in coefficient extraction the unknown contribution \(rD_n\varepsilon^n\) is omitted.  Every coefficient belongs to
\begin{equation}\label{eq:barnes-ring}
 D_n(h)\in\Q[c,h^{-1}],
 \qquad D_n(h)=\BigO_n(1)
 \quad(h\to+\infty).
\end{equation}
\end{theorem}

\begin{proof}
The coefficient of \(D_n\varepsilon^n\) in \cref{eq:barnes-rec-equation} is \(rD_n=(h/2)D_n\); all other contributions depend only on earlier coefficients.  This gives a triangular recurrence and proves uniqueness.  Since the input coefficients \(b_k\) are rational and \(a=c/2\), induction gives \(D_n\in\Q[c,h^{-1}]\).  Inspection of the recurrence after division by \(h/2\) shows that the coefficients remain bounded as \(h\to\infty\).  The first omitted displacement is \(D_{N+1}T^{-N-1}\).  Its logarithmic forward residual has size \(O(hT^{-N})\), because the derivative is asymptotic to \(Th/2\); division by that same derivative scale yields an inverse error \(O(T^{-N-1})\) for fixed \(N\).  The standard Barnes remainder estimate supplies the analytic input.
\end{proof}

\subsection{Explicit Barnes coefficients}

The first six coefficients are most readable when grouped by powers of
\(c=\log(2\pi)\):
\begin{align}
 D_0(h)&=-\frac{c}{h},\label{eq:D0}\\[0.2em]
 D_1(h)&=\frac1{12}+\frac1{6h}
 +c^2\left(\frac1{2h^2}-\frac1{h^3}\right),\label{eq:D1}\\[0.2em]
 D_2(h)&=\frac{c}{3h^3}
 +c^3\left(\frac{4}{3h^4}-\frac{2}{h^5}\right),\label{eq:D2}\\[0.2em]
 D_3(h)&=-\frac1{288}+\frac1{720h}-\frac1{72h^2}-\frac1{36h^3}\notag\\
 &\quad+c^2\left(-\frac1{24h^2}-\frac1{6h^3}-\frac1{6h^4}+\frac1{h^5}\right)\notag\\
 &\quad+c^4\left(-\frac1{8h^4}-\frac1{6h^5}
 +\frac{25}{6h^6}-\frac5{h^7}\right),\label{eq:D3}\\[0.2em]
 D_4(h)&=c\left(\frac{11}{360h^2}+\frac{11}{360h^3}
 -\frac1{18h^4}-\frac1{6h^5}\right)\notag\\
 &\quad+c^3\left(-\frac1{9h^4}-\frac7{9h^5}
 -\frac{10}{9h^6}+\frac{10}{3h^7}\right)\notag\\
 &\quad+c^5\left(-\frac4{5h^6}-\frac4{3h^7}
 +\frac{14}{h^8}-\frac{14}{h^9}\right),\label{eq:D4}\\[0.2em]
 D_5(h)&=\frac1{3456}-\frac{167}{45360h}-\frac1{240h^2}
 -\frac1{2160h^3}+\frac5{648h^4}+\frac1{108h^5}\notag\\
 &\quad+c^2\left(\frac1{192h^2}+\frac{17}{360h^3}
 +\frac1{6h^4}+\frac{31}{120h^5}
 -\frac5{36h^6}-\frac5{6h^7}\right)\notag\\
 &\quad+c^4\left(\frac1{32h^4}+\frac3{16h^5}+\frac1{24h^6}
 -\frac{55}{18h^7}-\frac{35}{6h^8}+\frac{35}{3h^9}\right)\notag\\
 &\quad+c^6\left(\frac1{16h^6}+\frac9{40h^7}-\frac{77}{20h^8}
 -\frac{70}{9h^9}+\frac{49}{h^{10}}-\frac{42}{h^{11}}\right).
 \label{eq:D5}
\end{align}

Consequently,
\begin{equation}\label{eq:barnes-six-sector}
\begin{split}
 \Barinv(y)=1+T+D_0(h)+\frac{D_1(h)}{T}
 +\frac{D_2(h)}{T^2}+\frac{D_3(h)}{T^3}
 +\frac{D_4(h)}{T^4}+\frac{D_5(h)}{T^5}
 +\BigO(T^{-6}).
\end{split}
\end{equation}

\subsection{The first Barnes corrections by hand}

The constant term of \cref{eq:barnes-rec-equation} gives
\[
 \frac h2D_0+\frac c2=0,
\]
so \(D_0=-c/h\).  At order \(\varepsilon\),
\[
 \frac h2D_1+\frac{\sigma}{2}D_0^2
 +\frac c2D_0-\frac{\sigma}{12}=0,
\]
which simplifies to \cref{eq:D1}.  At order \(\varepsilon^2\),
\[
 \frac h2D_2+\frac{D_0^3}{6}
 +\sigma D_0D_1+\frac c2D_1-\frac{D_0}{12}=0,
\]
and substitution gives \cref{eq:D2}.  The Bernoulli coefficient
\(b_1=-1/240\) first enters at order \(\varepsilon^3\), hence in \(D_3\).

\subsection{Elementary log--log expansion of the Barnes scale}\label{sec:barnes-loglog}

Put
\begin{equation}\label{eq:barnes-log-vars}
 A=\log(4Y),\qquad B=\log A,
 \qquad D=B+3.
\end{equation}
The exact dominant equation gives
\begin{equation}\label{eq:T-barnes-loglog}
\begin{split}
 T=\sqrt{\frac{4Y}{A}}\Bigg[&1+\frac{D}{2A}
 +\frac{D(3D-4)}{8A^2}
 +\frac{D(5D^2-16D+8)}{16A^3}\\
 &+\frac{D(105D^3-568D^2+720D-192)}{384A^4}\\
 &+\frac{D(189D^4-1488D^3+3296D^2-2304D+384)}{768A^5}
 +\BigO\!\left(\frac{B^6}{A^6}\right)\Bigg].
\end{split}
\end{equation}

\begin{proof}[Derivation]
Write \(T=\sqrt{4Y/A}\,F\).  Since
\(\log T=\tfrac12(A-B)+\log F\), the dominant equation is equivalent to
\begin{equation}\label{eq:Fbarnes}
 F^2\left(1-\frac{D}{A}+\frac{2\log F}{A}\right)=1.
\end{equation}
Set \(F=1+\sum_{n\ge1}f_nA^{-n}\) and equate powers of \(A^{-1}\).  The displayed coefficients follow successively.
\end{proof}

As in the Gamma case, the additive corrections \(D_nT^{-n}\) are beyond every relative order in the elementary log--log expansion of \(T\).  The leading Barnes displacement
\(-c/h\) is only logarithmically small in absolute size, but relative to
\(T\) it is exponentially small in \(w\).

\subsection{A resummable subsector}\label{sec:barnes-resum}

The limits of the first coefficients as \(h\to\infty\) are
\[
 D_0\to0,
 \quad D_1\to\frac1{12},
 \quad D_2\to0,
 \quad D_3\to-\frac1{288},
 \quad D_4\to0,
 \quad D_5\to\frac1{3456}.
\]
These are not unrelated constants.

\begin{proposition}[Square-root resummation of the leading constant subsector]\label{prop:barnes-resum}
For \(j\ge1\), the leading \(h^0\) part of \(D_{2j-1}(h)\) is
\begin{equation}\label{eq:barnes-limit-coeff}
 \lim_{h\to\infty}D_{2j-1}(h)
 =\binom{1/2}{j}6^{-j},
\end{equation}
and \(\lim_{h\to\infty}D_{2j}(h)=0\).  Therefore the entire leading constant subsector resums to
\begin{equation}\label{eq:barnes-sqrt-resum}
 T+\sum_{j\ge1}\binom{1/2}{j}6^{-j}T^{1-2j}
 =\sqrt{T^2+\frac16}.
\end{equation}
\end{proposition}

\begin{proof}
Divide \cref{eq:barnes-rec-equation} by \(h/2\) and let \(h\to\infty\) while treating \(\varepsilon\) formally.  Terms containing \(c/h\), the Bernoulli tail, and the logarithm without a factor of \(h\) vanish.  Since \(\sigma/h\to1/2\), the limiting equation is
\[
 \delta+\frac{\varepsilon}{2}\delta^2-\frac{\varepsilon}{12}=0.
\]
The solution of positive valuation is
\[
 \delta=\frac{\sqrt{1+\varepsilon^2/6}-1}{\varepsilon}
 =\sum_{j\ge1}\binom{1/2}{j}6^{-j}\varepsilon^{2j-1}.
\]
Because \(t=T+\delta\) and \(\varepsilon=T^{-1}\), this is exactly
\cref{eq:barnes-sqrt-resum}.
\end{proof}

\begin{remark}
The resummation does not replace the full expansion: the finite-\(h\) terms, especially \(-c/h\), are often numerically larger.  It does identify a coherent infinite subsector and suggests the improved internal scale
\(\widehat T=\sqrt{T^2+1/6}\) for very high-order symbolic work.
\end{remark}

\subsection{A direct Barnes residual certificate}

For
\[
 x_N=1+T+\sum_{n=0}^{N}D_n(h)T^{-n},
\]
the symbolic certificate is
\begin{equation}\label{eq:barnes-residual-cert}
 \log G(x_N)-\log y=\BigO(hT^{-N}),
\end{equation}
with the Barnes input series retained far enough to make the displayed order meaningful.  Since
\[
 \frac{\dd}{\dd x}\log G(x)
 \sim x(\log x-1)\sim \frac{Th}{2},
\]
the residual estimate gives
\[
 x_N-\Barinv(y)=\BigO(T^{-N-1}),
\]
in agreement with \cref{thm:barnes-all-orders}.  Again, the logarithmic residual is the stable quantity to evaluate numerically.

\subsection{Complex branch template}

On a logarithm sheet indexed by \(m\in\mathbb Z\), define
\[
 Y_m=\Log y+2\pi\ii m-\zeta'(-1).
\]
Choose a Lambert branch \(W_k\) and a square-root branch consistently:
\[
 w_{m,k}=W_k(4Y_m\e^{-3}),
 \qquad T_{m,k}=\left(\frac{4Y_m}{w_{m,k}}\right)^{1/2},
 \qquad h_{m,k}=w_{m,k}+1.
\]
Then the sectorial formal template is
\begin{equation}\label{eq:barnes-complex-template}
 x_{m,k}\sim1+T_{m,k}
 +\sum_{n\ge0}D_n(h_{m,k})T_{m,k}^{-n}.
\end{equation}
The square-root choice and the Lambert branch are linked by continuation; treating them as independent labels can double-count the same local branch.  As with Gamma, the formula is a local asymptotic chart rather than a global classification.

\section{Extensions of the general method}\label{sec:extensions}

\subsection{A reusable class of forward transseries}

The two special functions fit the common model
\begin{equation}\label{eq:general-forward-class}
 F(x)=A x^p(\log x-c)+C
 +\sum_{\lambda\in\Lambda}x^{p-\lambda}P_\lambda(\log x)
 +\mathcal R(x),
\end{equation}
where \(A\ne0\), \(p>0\), the support is locally finite, and
\(\mathcal R\) is a controlled remainder.  First absorb \(C\) into the right-hand side, then apply \cref{thm:lambert-normal} to the complete leading monomial.  The inverse has the form
\begin{equation}\label{eq:general-inverse-class}
 F^{-1}(Y)=X\left(1+
 \sum_{\mu\in\calS\setminus\{0\}}U_\mu(w)X^{-\mu}\right),
\end{equation}
where \(\calS\) is the additive semigroup generated by \(\Lambda\).  In absolute rather than relative form, the exponents are \(X^{1-\mu}\).

\begin{example}[Gamma]
Here \(A=1\), \(p=1\), \(c=1\), and after the half-shift and constant absorption the perturbation exponents are
\[
 \Lambda=\{2,4,6,\ldots\}.
\]
Thus \(\calS=2\N_0\), and the absolute inverse powers are
\(T^{1-2n}\).
\end{example}

\begin{example}[Barnes \(G\)]
Here \(A=\tfrac12\), \(p=2\), \(c=\tfrac32\), and after constant absorption
\[
 \Lambda=\{1,2,4,6,\ldots\}.
\]
The exponent \(1\) comes from the linear term, so the generated semigroup is all of \(\N_0\); every integer absolute correction power can occur.
\end{example}

\subsection{The full-leading-monomial rule}

A frequent failure mode is to normalize only by \(A x^p\), leaving
\(\log x-c\) in the perturbation.  That ``perturbation'' is not small, and the resulting near-identity recursion either fails or silently creates an extra logarithmic generator.  The correct dominant object is the full monomial
\[
 A x^p(\log x-c),
\]
including its logarithmic factor and constant shift.  Its exact inverse produces the correct slow coordinate \(w\) and the correct divisor \(w+1\).

A second failure mode is to neglect a useful translation of \(x\).  A shift can remove a large family of subleading powers before inversion.  The Gamma shift \(x=t+1/2\) is the canonical example.  In a general problem one should first search for a translation
\(x=t+\alpha\) that simplifies the coefficients of \(t^{p-1}\log t\),
\(t^{p-1}\), or the parity of the inverse-power tail.

\subsection{Support propagation and resonance}

The inverse support is not usually the original support \(\Lambda\), but its additive closure.  A perturbation containing exponents \(2\) and \(3\), for example, generally produces inverse exponents
\[
 2,3,4,5,6,\ldots
\]
through nonlinear products.  Local finiteness is therefore the right condition: it permits infinitely many exponents overall while guaranteeing finite coefficient arithmetic at every requested order.

The factor \(w+1\) is a unit for the large positive branch, but it can become small on complex branches near a Lambert branch point.  This is an inverse resonance.  The formal coefficient formulas still display what goes wrong---their poles accumulate powers of \((w+1)^{-1}\)---while the analytic residual-to-error bound simultaneously loses strength because the dominant derivative vanishes.

\subsection{Several logarithmic levels}

If the leading forward monomial is
\[
 A x^p(\log x)^q,
\]
with \(q\ne1\), one may still obtain a Lambert-type dominant inverse after a power substitution in special cases, or use a generalized Lambert function.  More generally, taking logarithms gives
\[
 p\log x+q\log\log x=\log(Y/A),
\]
which is a near-identity equation for \(p\log x\) with a second logarithmic level.  The same strategy applies:

\begin{enumerate}
\item choose the complete dominant transmonomial;
\item solve it exactly when a named inverse is available, or construct its own nested-log inverse;
\item express all remaining terms in a well-ordered small-monomial group;
\item revert the normalized displacement equation by triangular recursion or Newton iteration.
\end{enumerate}

The essential theorem is not tied to Lambert \(W\); it is an implicit-function theorem in a valued transseries field.  Lambert \(W\) is exceptionally useful because it packages the entire first logarithmic inversion into one exact analytic coordinate.

\subsection{Multiple Gamma functions}

Higher Barnes multiple Gamma functions have logarithmic asymptotics whose leading terms are polynomial multiples of \(x^r\log x\).  Once the dominant coefficient and normalization are fixed, their large inverse branches fall under the same scheme with \(p=r\).  The perturbation support then records lower polynomial powers, logarithms, and the Bernoulli tail.  The present Barnes \(G\) calculation is the first nontrivial member beyond Euler Gamma and provides a template for systematic higher-order work.

\subsection{Conditioning and practical stopping rules}

The asymptotic inverse condition number is encoded by
\[
 \frac{1}{F'(X)}\sim\frac{1}{A X^{p-1}(w+1)}.
\]
This suggests a practical stopping rule.  Given a candidate \(x_N\), compute the logarithmic residual \(r_N=F(x_N)-Y\) in the same logarithmic equation used for derivation.  Stop when
\[
 \frac{|r_N|}{|F'(x_N)|}
\]
is below the target absolute tolerance.  This approach automatically adapts to the branch, avoids overflow, and remains meaningful when the asymptotic series has reached its least term.

\section{Exponentially small sectors and Stokes inheritance}\label{sec:stokes}

The algebraic inverse transseries derived above is complete relative to fixed-order Poincar\'e input.  It does not include the exponentially improved sectors associated with the late terms of the Gamma and Barnes expansions.

\begin{proposition}[First inherited exponential correction]\label{prop:stokes-inheritance}
Suppose a logarithmic forward equation has been split as
\[
 F(x)=F_{\mathrm{alg}}(x)+E(x)=Y,
\]
where \(x_{\mathrm{alg}}\) solves \(F_{\mathrm{alg}}(x)=Y\) to the desired algebraic-transseries order, and \(E(x)\) is exponentially smaller.  If
\(F_{\mathrm{alg}}'(x_{\mathrm{alg}})\ne0\), then the first exponential correction to the inverse is
\begin{equation}\label{eq:first-exp-inverse}
 \Delta x
 =-\frac{E(x_{\mathrm{alg}})}{F_{\mathrm{alg}}'(x_{\mathrm{alg}})}
 +\BigO\!\left(
 \frac{E(x_{\mathrm{alg}})E'(x_{\mathrm{alg}})}{F_{\mathrm{alg}}'(x_{\mathrm{alg}})^2}
 +\frac{E(x_{\mathrm{alg}})^2F_{\mathrm{alg}}''(x_{\mathrm{alg}})}{F_{\mathrm{alg}}'(x_{\mathrm{alg}})^3}
 \right).
\end{equation}
\end{proposition}

\begin{proof}
Insert \(x=x_{\mathrm{alg}}+\Delta x\), Taylor-expand both terms, and solve the linearized equation.  The displayed remainder records the two quadratic sources: variation of \(E\) and curvature of \(F_{\mathrm{alg}}\).
\end{proof}

For inverse Gamma, a forward contribution of size \(E(T)\) is divided by
\(\psi(T+1/2)\sim\log T=s\).  For inverse Barnes \(G\), it is divided by
\(T(\log T-1)\sim Th/2\).  Thus the inverse Stokes multipliers are the forward multipliers transported through the derivative.  The Stokes locations are inherited from the forward special function; inversion changes amplitudes and branch charts but does not create new exponential actions at first order.

Nemes' exponentially improved expansion for \(\log G(t+1)\) \cite{NemesBarnes} exhibits terminant-weighted terms with actions \(\exp(\pm2\pi\ii k t)\), whose multipliers change smoothly across \(\arg t=\pm\pi/2\).  Substitution of the algebraic inverse transseries into those actions, followed by \cref{eq:first-exp-inverse}, produces a transseries beyond the scope of the fixed-order formulas in \cref{sec:barnes}.  A complete resurgent treatment would also have to track how complex inverse branches cross the forward Stokes geometry.

\begin{warningbox}[Do not append exponentials by pattern matching]
The exponentials, their terminant factors, and their Stokes multipliers are sector-dependent.  A correct exponentially improved inverse must start from a correct exponentially improved forward expansion on the same branch and in the same sector.  Merely adding \(\e^{-2\pi T}\)-type terms to the real formula is not a branch-complete construction.
\end{warningbox}

\section{Numerical validation}\label{sec:numerics}

The tables below avoid numerical inversion altogether.  A test value \(x\) is chosen, the exact forward value \(y=\Gamma(x)\) or \(y=G(x)\) is computed at high precision, and the asymptotic inverse formula is evaluated at that \(y\).  The signed error is then approximation minus the known \(x\).  This procedure tests the complete chain from the forward special function to the inverse transseries.

\subsection{Gamma}

Let \(\mathcal G_N\) denote \cref{eq:gamma-all-orders} through \(C_N\), with
\(\mathcal G_0=\tfrac12+T\).  The signed errors are:

\begin{center}
\small
\begin{tabular}{r@{\qquad}rrrrrr}
\toprule
\(x\) & \(\mathcal G_0-x\) & \(\mathcal G_1-x\) & \(\mathcal G_2-x\) & \(\mathcal G_3-x\) & \(\mathcal G_4-x\) & \(\mathcal G_5-x\)\\
\midrule
5
& \(-6.1414\,10^{-3}\) & \(2.8694\,10^{-5}\) & \(-4.1862\,10^{-7}\)
& \(1.3452\,10^{-8}\) & \(-8.2694\,10^{-10}\) & \(8.4558\,10^{-11}\)\\
10
& \(-1.9470\,10^{-3}\) & \(1.7430\,10^{-6}\) & \(-5.7560\,10^{-9}\)
& \(4.4485\,10^{-11}\) & \(-6.5578\,10^{-13}\) & \(1.5870\,10^{-14}\)\\
20
& \(-7.1924\,10^{-4}\) & \(1.4126\,10^{-7}\) & \(-1.1183\,10^{-10}\)
& \(2.1113\,10^{-13}\) & \(-7.5612\,10^{-16}\) & \(4.4191\,10^{-18}\)\\
50
& \(-2.1572\,10^{-4}\) & \(6.1957\,10^{-9}\) & \(-7.6816\,10^{-13}\)
& \(2.2927\,10^{-16}\) & \(-1.2910\,10^{-19}\) & \(1.1809\,10^{-22}\)\\
100
& \(-9.1031\,10^{-5}\) & \(6.2869\,10^{-10}\) & \(-1.9388\,10^{-14}\)
& \(1.4440\,10^{-18}\) & \(-2.0231\,10^{-22}\) & \(4.5949\,10^{-26}\)\\
\bottomrule
\end{tabular}
\end{center}

Here \(a\,10^b\) means \(a\times10^b\).  Even at \(x=5\), which is not especially asymptotic, the sectors decrease rapidly through the displayed order.  At \(x=20\), four correction sectors reduce the error from approximately \(7.2\times10^{-4}\) to approximately \(7.6\times10^{-16}\).

\subsection{Barnes \texorpdfstring{\(G\)}{G}}

Let \(\mathcal B_{-1}=1+T\), and let \(\mathcal B_N\) include
\(D_0,\ldots,D_N\).  The signed errors are:

\begin{center}
\scriptsize
\setlength{\tabcolsep}{4.3pt}
\begin{tabular}{r@{\quad}rrrrrrrr}
\toprule
\(x\) & \(\mathcal B_{-1}-x\) & \(\mathcal B_0-x\) & \(\mathcal B_1-x\)
& \(\mathcal B_2-x\) & \(\mathcal B_3-x\) & \(\mathcal B_4-x\) & \(\mathcal B_5-x\)\\
\midrule
5
& \(1.3817\) & \(3.6331\,10^{-2}\) & \(-3.5643\,10^{-3}\)
& \(-3.3157\,10^{-3}\) & \(-1.3188\,10^{-3}\) & \(-3.9978\,10^{-4}\) & \(-9.4554\,10^{-5}\)\\
10
& \(7.0002\,10^{-1}\) & \(-2.2345\,10^{-2}\) & \(-1.2469\,10^{-3}\)
& \(1.0080\,10^{-5}\) & \(6.0183\,10^{-6}\) & \(2.2382\,10^{-7}\) & \(-2.8004\,10^{-8}\)\\
20
& \(4.5759\,10^{-1}\) & \(-9.2943\,10^{-3}\) & \(-8.0035\,10^{-5}\)
& \(2.8480\,10^{-6}\) & \(4.8454\,10^{-8}\) & \(-2.0959\,10^{-9}\) & \(-3.5723\,10^{-11}\)\\
50
& \(3.1413\,10^{-1}\) & \(-2.9440\,10^{-3}\) & \(-3.4332\,10^{-6}\)
& \(9.5141\,10^{-8}\) & \(-7.7542\,10^{-12}\) & \(-6.5727\,10^{-12}\) & \(5.0468\,10^{-14}\)\\
100
& \(2.5411\,10^{-1}\) & \(-1.3107\,10^{-3}\) & \(-4.0663\,10^{-7}\)
& \(8.4132\,10^{-9}\) & \(-5.5679\,10^{-12}\) & \(-9.5506\,10^{-14}\) & \(6.4482\,10^{-16}\)\\
\bottomrule
\end{tabular}
\end{center}

The Barnes expansion is less spectacular at small \(x\), as expected from its denser support and larger linear displacement, but the predicted order-by-order improvement becomes clear by \(x=10\) and is strong by \(x=20\).  As with every divergent asymptotic expansion, one should stop near the least term rather than assume monotone improvement forever; the slight comparison between \(\mathcal B_3\) and \(\mathcal B_4\) at \(x=50\) illustrates that finite-precision cancellation and asymptotic ordering can interact.

\subsection{What the tests do and do not prove}

The computations strongly check the algebra, signs, shifts, and branch normalizations.  They are not substitutes for the analytic remainder argument.  Conversely, a theorem proving only the order of the remainder would not catch a transcription error in a large rational coefficient.  Formal residual cancellation, high-precision substitution, and analytic remainder bounds are complementary checks.

\section{A reproducible symbolic workflow}\label{sec:code}

The following compact SymPy code regenerates the coefficients from the two defining recurrence equations using exact arithmetic.  It is intentionally based on residual coefficient extraction rather than on a hard-coded coefficient list.

\begin{lstlisting}[style=pythonstyle,caption={Exact symbolic generation of the inverse coefficients.}]
import sympy as sp


def gamma_coefficients(order: int):
    """Return C_1(s),...,C_order(s) for inverse Gamma."""
    eps, s = sp.symbols("eps s")
    C = sp.symbols(f"C1:{order + 1}")
    u = sum(C[n - 1] * eps**n for n in range(1, order + 1))

    # B_{2n}(1/2)/(2n(2n-1))
    a = {
        n: sp.bernoulli(2*n, sp.Rational(1, 2))
           / (sp.Integer(2*n) * (2*n - 1))
        for n in range(1, order + 1)
    }

    residual = s*u
    for m in range(2, order + 2):
        residual += (-1)**m * u**m / (sp.Integer(m) * (m - 1))
    for n in range(1, order + 1):
        residual += a[n] * eps**n * (1 + u)**(-(2*n - 1))

    residual = sp.series(residual, eps, 0, order + 1).removeO()
    solved = {}
    for n in range(1, order + 1):
        equation = sp.expand(residual.subs(solved)).coeff(eps, n)
        solved[C[n - 1]] = sp.factor(sp.solve(equation, C[n - 1])[0])
    return [solved[z] for z in C]


def barnes_coefficients(order: int):
    """Return D_0(h),...,D_order(h) for inverse Barnes G."""
    eps, h, c = sp.symbols("eps h c")
    D = sp.symbols(f"D0:{order + 1}")
    delta = sum(D[n] * eps**n for n in range(order + 1))

    r = h / 2
    sigma = 1 + h / 2
    a = c / 2

    residual = r*delta + a
    residual += eps * (
        sigma*delta**2/2 + a*delta - sigma/sp.Integer(12)
    )
    for m in range(3, order + 3):
        residual += (
            (-1)**(m - 3) * delta**m * eps**(m - 1)
            / (sp.Integer(m) * (m - 1) * (m - 2))
        )
    residual -= eps * sp.log(1 + eps*delta) / 12

    for k in range(1, order + 1):
        b_k = sp.bernoulli(2*k + 2) / (sp.Integer(4)*k*(k + 1))
        residual += (
            b_k * eps**(2*k + 1) * (1 + eps*delta)**(-2*k)
        )

    residual = sp.series(residual, eps, 0, order + 1).removeO()
    solved = {}
    for n in range(order + 1):
        equation = sp.expand(residual.subs(solved)).coeff(eps, n)
        solved[D[n]] = sp.factor(sp.solve(equation, D[n])[0])
    return [solved[z] for z in D]


if __name__ == "__main__":
    print("Gamma coefficients:")
    for n, value in enumerate(gamma_coefficients(5), start=1):
        print(f"C_{n} =", value)

    print("\nBarnes coefficients:")
    for n, value in enumerate(barnes_coefficients(5)):
        print(f"D_{n} =", sp.collect(sp.expand(value), sp.symbols("c")))
\end{lstlisting}

\subsection{Independent verification protocol}

For a trustworthy high-order computation, use at least the following checks.

\begin{enumerate}
\item Generate coefficients by the triangular recurrence.
\item Generate the same truncation by formal Newton iteration.
\item Substitute the result into the defining asymptotic equation and verify the first nonzero residual block exactly.
\item Evaluate the inverse approximation at \(y=\Gamma(x)\) or \(y=G(x)\) for several exact test values \(x\) using high-precision arithmetic.
\item Increase both working precision and the number of retained forward asymptotic terms.  Stable digits should not depend on either once both are adequate.
\end{enumerate}

The residual computation should be implemented independently of the coefficient generator whenever possible.  A routine that verifies its own output using the same erroneous composition rule can pass a circular test.

\section{Further deductions and research directions}\label{sec:research}

\subsection{Closed descriptions of coefficient families}

The recurrences prove existence and compute arbitrary order, but they do not yet give a satisfying closed combinatorial formula for \(C_n(s)\) or \(D_n(h)\).  The Gamma \cref{eq:gamma-rec-equation} combines the coefficients of
\((1+u)\log(1+u)\), binomial powers, and Bernoulli values at \(1/2\).  This strongly suggests expressions in partial Bell polynomials decorated by Bernoulli numbers.  A useful goal is a formula that makes the observed denominator pattern
\[
 C_n(s)=\frac{P_n(s)}{s^{2n-1}}
\]
manifest, together with degree, sign, and coefficient properties of \(P_n\).

For Barnes \(G\), the grouping by powers of \(c=\log(2\pi)\) reveals a parity rule:
\(D_n\) contains only powers \(c^j\) with \(j\equiv n+1\pmod2\); when \(n\) is odd, this includes the constant power \(c^0\).  The recurrence strongly suggests an elementary grading proof and a concise combinatorial formulation.

\subsection{Optimal truncation of the inverse}

The fixed-order theorems do not identify the least term as the order grows with \(T\).  Because the forward Bernoulli coefficients grow factorially, the inverse coefficient families are expected to inherit factorial-over-power growth.  Determining their late-term ansatz would reveal the optimal truncation index and the precise exponential remainder of the inverse.  The inherited-exponential formula \cref{prop:stokes-inheritance} gives the first structural clue, but a complete analysis should derive the inverse singulant and Stokes multiplier directly.

\subsection{Uniformity near branch degeneration}

The expansions are nonuniform when \(w+1\) is small.  Near \(w=-1\), the dominant map has a quadratic turning point, and the correct local inverse scale is a square root rather than a regular power series in the residual.  A uniform approximation should match a branch-point expansion in \((w+1)\) to the large-\(w\) transseries developed here.  This problem is relevant to complex inverse branches even though the positive-real large-argument branch stays far from degeneration.

\subsection{Higher Barnes hierarchy}

Applying the method to Barnes multiple Gamma functions should produce a family of Lambert scales with \(p=3,4,\ldots\), followed by increasingly rich support semigroups.  The main technical tasks are:

\begin{enumerate}
\item choose translations that simplify the generalized Bernoulli-polynomial tails;
\item identify constants to absorb into the logarithmic right-hand side;
\item compute and classify the inverse coefficient rings;
\item determine the higher-order analogues of the Barnes square-root resummation in \cref{prop:barnes-resum}.
\end{enumerate}

\subsection{Computer algebra and formal verification}

The triangular recurrence is well suited to a theorem prover because every requested coefficient depends on finitely many earlier blocks.  A formal development can separate:

\begin{itemize}
\item the algebra of locally finite supports;
\item the formal implicit-function theorem and Newton precision doubling;
\item exact rational coefficient certificates;
\item analytic lemmas importing the special-function remainder estimates;
\item interval or ball arithmetic for numerical residual bounds.
\end{itemize}

This separation mirrors the distinction between a formal transseries identity and an asymptotic statement about an analytic realization.

\section{Conclusion}

The apparent difficulty of inverting Gamma and Barnes \(G\) at large values comes from trying to reverse all scales at once.  After taking logarithms and choosing the correct shift, the full dominant logarithmic monomial is inverted exactly by Lambert \(W\).  The rest is a regular, triangular near-identity problem in a fast exponential scale with slow rational coefficients.

For Gamma, the half-shift reduces the entire correction support to even relative powers, producing only odd absolute inverse powers.  For Barnes \(G\), the linear Stirling term creates a logarithmic displacement and fills the integer support; nevertheless, every coefficient remains an explicit polynomial in \(\log(2\pi)\) and \((w+1)^{-1}\).  Both expansions admit elementary log--log re-expansions, but the exact \(W\)-coordinates preserve the hierarchy more transparently and perform better numerically.

The general theorem applies well beyond these examples.  Under the support, nondegeneracy, and analytic-remainder hypotheses stated above, a forward asymptotic expansion built from
\(A x^p(\log x-c)\) and a locally finite lower transseries can be inverted by the same sequence: exact Lambert normalization, support closure, triangular or Newton reversion, residual certification, and analytic error transfer.  This provides a practical bridge between formal transseries algebra, special-function asymptotics, and reliable computation.

\appendix

\section{Derivation of the normalized Barnes equation}\label{app:barnes-derivation}

For reference, we derive \cref{eq:barnes-rec-equation} without skipping algebra.  Put
\[
 F_0(t)=\frac12t^2\log t-\frac34t^2,
 \qquad t=T(1+v),\qquad v=\varepsilon\delta.
\]
Since \(F_0(T)=Y\) and \(\log T=\sigma\),
\begin{align*}
 \frac{F_0(T(1+v))-F_0(T)}{T}
 &=\frac1{2\varepsilon}
 \left[(1+v)^2\left(\sigma-\frac32+\log(1+v)\right)
 -\left(\sigma-\frac32\right)\right].
\end{align*}
Now
\begin{equation}\label{eq:oneplus-square-log}
 (1+v)^2\log(1+v)
 =v+\frac32v^2+
 \sum_{m\ge3}\frac{2(-1)^{m-3}}{m(m-1)(m-2)}v^m.
\end{equation}
Using \(r=\sigma-1=h/2\) gives
\begin{equation}\label{eq:F0-barnes-diff}
 \frac{F_0(T+\delta)-F_0(T)}{T}
 =r\delta+\frac{\sigma}{2}\varepsilon\delta^2
 +\sum_{m\ge3}\frac{(-1)^{m-3}}{m(m-1)(m-2)}
 \varepsilon^{m-1}\delta^m.
\end{equation}
The remaining pieces of \cref{eq:barnes-simplified}, after division by \(T\), are
\begin{align*}
 \frac{ct/2}{T}&=\frac c2+\frac c2\varepsilon\delta,\\
 -\frac{\log t}{12T}
 &=-\frac{\sigma}{12}\varepsilon
 -\frac{\varepsilon}{12}\log(1+\varepsilon\delta),\\
 \frac{b_k t^{-2k}}{T}
 &=b_k\varepsilon^{2k+1}(1+\varepsilon\delta)^{-2k}.
\end{align*}
Adding these expressions proves \cref{eq:barnes-rec-equation}.

\section{The first Gamma coefficient equations}\label{app:gamma-equations}

Writing \(u=\sum_{n\ge1}C_n\varepsilon^n\), the first five triangular equations before simplification are
\begin{align*}
0={}&sC_1+a_1,\\
0={}&sC_2+\frac12C_1^2-a_1C_1+a_2,\\
0={}&sC_3+C_1C_2-\frac16C_1^3
 +a_1(C_1^2-C_2)-3a_2C_1+a_3,\\
0={}&sC_4+C_1C_3+\frac12C_2^2-\frac12C_1^2C_2
 +\frac1{12}C_1^4\\
&\quad+a_1(-C_3+2C_1C_2-C_1^3)
 +a_2(-3C_2+6C_1^2)-5a_3C_1+a_4,\\
0={}&sC_5+\text{a finite polynomial in }C_1,C_2,C_3,C_4
 \text{ and }a_1,\ldots,a_5.
\end{align*}
The fifth polynomial is best generated mechanically; its simplified solution is \cref{eq:C5}.  These equations also show directly that each new coefficient occurs linearly with divisor \(s\).

\section{Elementary-log coefficient recurrences}\label{app:log-recurrences}

The Gamma elementary-log coefficients \(f_n(D)\) are defined by
\[
 F=1+\sum_{n\ge1}f_n(D)A^{-n},
 \qquad
 F\left(1-\frac{D}{A}+\frac{\log F}{A}\right)=1.
\]
The first six are
\begin{align*}
 f_1&=D,\\
 f_2&=D(D-1),\\
 f_3&=\frac12D(D-2)(2D-1),\\
 f_4&=\frac16D(6D^3-26D^2+27D-6),\\
 f_5&=\frac1{12}D(12D^4-77D^3+142D^2-84D+12),\\
 f_6&=\frac1{30}D(30D^5-261D^4+725D^3-775D^2+300D-30).
\end{align*}

For Barnes \(G\), define
\[
 F=1+\sum_{n\ge1}g_n(D)A^{-n},
 \qquad
 F^2\left(1-\frac{D}{A}+\frac{2\log F}{A}\right)=1.
\]
Then
\begin{align*}
 g_1&=\frac D2,\\
 g_2&=\frac18D(3D-4),\\
 g_3&=\frac1{16}D(5D^2-16D+8),\\
 g_4&=\frac1{384}D(105D^3-568D^2+720D-192),\\
 g_5&=\frac1{768}D(189D^4-1488D^3+3296D^2-2304D+384),\\
 g_6&=\frac1{15360}D(3465D^5-36516D^4+120400D^3
 -150400D^2+67200D-7680).
\end{align*}

\section{Notation summary}\label{app:notation}

\begin{longtable}{@{}p{0.19\textwidth}p{0.73\textwidth}@{}}
\toprule
Symbol & Meaning\\
\midrule
\endhead
\(y\) & Large argument of the inverse special function.\\
\(x\) & Inverse value: \(\Gamma(x)=y\) or \(G(x)=y\).\\
\(t=x-1/2\) & Symmetry-adapted Gamma inverse variable.\\
\(t=x-1\) & Barnes inverse variable matching \(G(t+1)\).\\
\(W_k\) & Branch \(k\) of Lambert's \(W\)-function, satisfying \(W_k(z)e^{W_k(z)}=z\).\\
\(M\) & Gamma logarithmic right-hand side \(\log(y/\sqrt{2\pi})\).\\
\(Y\) & Barnes logarithmic right-hand side \(\log y-\zeta'(-1)\).\\
\(T\) & Exact inverse of the dominant logarithmic monomial.\\
\(s=w+1\) & Gamma normalized derivative and \(\log T\).\\
\(h=w+1\) & Barnes normalized derivative parameter; \(\log T=1+h/2\).\\
\(C_n(s)\) & Gamma fast-sector coefficients.\\
\(D_n(h)\) & Barnes fast-sector coefficients.\\
\(a_n\) & Shifted Stirling coefficients \(B_{2n}(1/2)/(2n(2n-1))\).\\
\(b_k\) & Barnes tail coefficients \(B_{2k+2}/(4k(k+1))\).\\
\(\calS\) & Additive semigroup generated by the perturbation exponents.\\
\bottomrule
\end{longtable}

\begingroup
\sloppy
\begin{thebibliography}{99}

\bibitem{DLMF}
NIST Digital Library of Mathematical Functions,
Chapters 4 and 5, especially \S\S4.13, 5.11, and 5.17,
\url{https://dlmf.nist.gov/}, accessed 3 September 2026.

\bibitem{CorlessW}
R.~M. Corless, G.~H. Gonnet, D.~E.~G. Hare, D.~J. Jeffrey, and D.~E. Knuth,
``On the Lambert \(W\) function,''
\emph{Advances in Computational Mathematics} \textbf{5} (1996), 329--359.

\bibitem{BorweinCorless}
J.~M. Borwein and R.~M. Corless,
``Gamma and factorial in the Monthly,''
\emph{American Mathematical Monthly} \textbf{125} (2018), 400--424,
\url{https://doi.org/10.1080/00029890.2018.1420983}.

\bibitem{Uchiyama}
M. Uchiyama,
``The principal inverse of the gamma function,''
\emph{Proceedings of the American Mathematical Society} \textbf{140} (2012),
1343--1348,
\url{https://doi.org/10.1090/S0002-9939-2011-11023-2}.

\bibitem{Pedersen}
H.~L. Pedersen,
``Inverses of gamma functions,''
\emph{Constructive Approximation} \textbf{41} (2015), 251--267,
\url{https://doi.org/10.1007/s00365-014-9239-1}.

\bibitem{FolitseJeffreyCorless}
K.~A. Folitse, D.~J. Jeffrey, and R.~M. Corless,
``Properties and computation of the functional inverse of Gamma,''
in \emph{2017 19th International Symposium on Symbolic and Numeric Algorithms for Scientific Computing}, IEEE, 2017, 63--66,
\url{https://doi.org/10.1109/SYNASC.2017.00020}.

\bibitem{JeffreyWatt}
D.~J. Jeffrey and S.~M. Watt,
``The inverse of the complex Gamma function,''
\emph{25th International Symposium on Symbolic and Numeric Algorithms for Scientific Computing}, 2023,
\url{https://doi.org/10.1109/SYNASC61333.2023.00010}; preprint
\url{https://arxiv.org/abs/2311.16583}.

\bibitem{Barnes1900}
E.~W. Barnes,
``The theory of the \(G\)-function,''
\emph{Quarterly Journal of Pure and Applied Mathematics} \textbf{31} (1900),
264--314.

\bibitem{FerreiraLopez}
C. Ferreira and J.~L. L\'opez,
``An asymptotic expansion of the double gamma function,''
\emph{Journal of Approximation Theory} \textbf{111} (2001), 298--314.

\bibitem{NemesBarnes}
G. Nemes,
``Error bounds and exponential improvement for the asymptotic expansion of the Barnes \(G\)-function,''
\emph{Proceedings of the Royal Society A} \textbf{470} (2014), 20140534,
\url{https://doi.org/10.1098/rspa.2014.0534}.

\bibitem{SalvyShackell}
B. Salvy and J. Shackell,
``Symbolic asymptotics: multiseries of inverse functions,''
\emph{Journal of Symbolic Computation} \textbf{27} (1999), 543--563,
\url{https://doi.org/10.1006/jsco.1999.0281}.

\bibitem{vanderHoeven}
J. van der Hoeven,
\emph{Transseries and Real Differential Algebra},
Lecture Notes in Mathematics 1888, Springer, 2006.

\bibitem{Olver}
F.~W.~J. Olver,
\emph{Asymptotics and Special Functions},
A K Peters, 1997 reprint.

\bibitem{ProveItPLT}
V. Reshetnikov,
``Polynomial--Logarithmic Transseries: Algebra, Composition, Series Reversal, and the Lambert \(W\) Archetype,''
ProveIt research draft, 2 September 2026,
\url{https://github.com/VladimirReshetnikov/ProveIt/tree/main/Analysis/FabiusFunction/docs/semi-formalized-research-frontiers/drafts/series-and-transseries}.

\end{thebibliography}
\endgroup

\end{document}
